\scanpage{20}
\uline{Cohomologie des quadriques}

$Q^n$ quadrique non singulière de dim $n$ sur $k$ sép. clos.

$Q^{n-1}$ section à l'infini.

\[
U^n=AQ^n=Q^n-Q^{n-1}
\qquad\text{\unsure{supposons|pposons|pour} $A$ coefficient annulé par $m$}
\]
\[
\begin{aligned}
H_c^i(U^n,A)&\simeq
\begin{cases}
0&\text{si }i\ne n,2n\\
A\otimes V'_n\deleted{\unsure{(i)|0}}&\text{si }i=n\\
A\otimes\mu_m(-n)&\text{si }i=2n
\end{cases}\quad\inserted{\text{si }n\ne0}\\
V'_n&=\check V_n\otimes\mu_m(-n)=\begin{cases}
\mu_m(-\frac n2)&n\text{ pair}\\
\mu_m(\frac{\unsure{-(n+1)|-n+1}}2)&n\text{ imp.}
\end{cases}\\
V'_n&=V_n^{\unsure{*|+|\vee}}=\mu_m(-\tfrac n2)\qquad n\text{ pair}\\
V'_n&\unsure{\simeq|=} V_n\otimes\mu_m(1)\qquad n\text{ impair}
\end{aligned}
\]
% slanted
\margin{Si $n$ impair, ces isom. ne sont déterminés qu'au signe près ....}
\[
\Updownarrow
\]
dualement
\[
\begin{aligned}
H^i(U^n,A)&\simeq\begin{cases}
0 &\text{si }i\ne0,n\\
A\otimes V_n &\text{si }i=0\\
A &\text{si }i=n.
\end{cases}\quad\inserted{\text{si }n\ne0\ \substack{\nearrow\\\nearrow}}\\
V_n&\simeq\mu_m(-\tfrac n2)\quad\text{si $n$ pair,}\\
V_n&\simeq\mu_m(\tfrac{\unsure{-n+1|-(n+1)|n\pm1}}2)\quad\text{si $n$ impair}
\end{aligned}
\]
$H_c^i(U^n,A)\xrightarrow{\varphi}H_c^i(U^n,A)$ est \uline{nul} \uline{sauf} si $n$ pair, $i=n$, dans ce cas \unsure{devient} un \unsure{isomorphisme} \inserted{exc. 2}.
\[
\deleted{\unsure{H^i(Q^n)\simeq H^i(Q^n)|H^i(Q)\simeq H^i(Q)}\longrightarrow}
\]
Il y a \inserted{deux} \unsure{un gén.|un génér.} $l$ de $H_c^i(U^n,A)$ tels que $\langle l,\varphi(l)\rangle=2$.

\diagram{\begin{tikzcd}[ampersand replacement=\&,column sep=tiny]
H^{i-1}(U^n)\arrow[r,dash] \& H^{i-2}(Q^{n-1})(-1)\arrow[r] \& H^i(Q^n)\arrow[r] \& H^i(\deleted{\unsure{Q}}U^n)\arrow[r,dash] \& H^{i-1}(Q^{n-1})(-1)\arrow[r,dash] \& \cdots
\end{tikzcd}}
\[
H^{i-2}(Q^{n-1})(-1)\xrightarrow{\alpha_*}H^i(Q^n)\quad\text{est}\quad
\left|\begin{aligned}
&\text{surjectif si }i\ne0,n\\
&\text{injectif si }i\ne1,n+1
\end{aligned}\right.
\quad\left|\begin{aligned}&\text{bijectif si}\\&i\ne0,\deleted{\unsure{1}}n,n+1\end{aligned}\right.
\]

Que se passe-t-il \unsure{pr|pour} $i=0,1,n,n+1$ ??

Pour $i=0,1$ $\alpha_*$ est \uline{nul} [$i=1$, \deleted{bij} $\alpha$ bijectif, car les deux mbs nuls)

Pour $i=n,n+1$ on a la \deleted{\unsure{suite}} \unsure{exacte} (si $\boxed{n\ne1}$ \unsure{donc} $n-1\ne0$]
\diagram{\begin{tikzcd}[ampersand replacement=\&,column sep=tiny]
0\arrow[r] \& H^{n-2}(Q^{n-1})\deleted{(-1)}\arrow[r] \& H^n(\deleted{\unsure{U}}Q^n)\arrow[r] \& H^n(U^n)\arrow[r,dash] \& H^{n-1}(Q^{n-1})(-1)\arrow[r,dash] \& H^{n+1}(Q^n)\arrow[r,dash] \& 0
\end{tikzcd}}

1) $n$ \uline{pair} $>0$, la suite exacte se réduit à
\diagram{\begin{tikzcd}[ampersand replacement=\&,column sep=small]
0\arrow[r] \& H^{n-2}(Q^{n-1})\deleted{(-1)}\arrow[r] \& H^n(Q^n)\arrow[r] \& H^n(U^n)\arrow[r] \& 0
\end{tikzcd}}
donc $H^n(Q^n)$ est extension de $H^n(U^n)$ par $H^{n-2}(Q^{n-1})(-1)$ ...
% archive sheet number: 19
